% ============================================================================
% Fichier  : partie2/P02_C06_recevabilite.tex : FORN (LPI342)
% Sources : ISO/IEC 27037, 27042 ; ACPO ; principes de la preuve.
% ============================================================================

\chapter{Recevabilité et Admissibilité}
\label{chap:recevabilite}
\chaptermark{Recevabilité et admissibilité}

\epigraph{%
    « Une vérité que l'on ne peut pas défendre\\
    n'entre pas dans un prétoire. »%
}{Principe du cours}

\niveauChapitre{Intermédiaire}{Chapitres~\ref{chap:droit} et~\ref{chap:custody}.}

\begin{boxobjectifs}
\begin{itemize}
    \item Définir la recevabilité d'une preuve numérique.
    \item Connaître ses conditions (licéité, intégrité, authenticité, traçabilité).
    \item Comprendre les quatre principes de l'ACPO.
    \item Relier la recevabilité au Trilemme CRO.
    \item Reconnaître ce qui rend une preuve irrecevable.
\end{itemize}
\end{boxobjectifs}

\bigskip

% ============================================================================
\section*{Question directrice : qu'est-ce qui rend une preuve numérique recevable, ou irrecevable, devant la justice ?}
% ============================================================================

Une preuve peut être exacte et pourtant écartée~: obtenue sans titre,
altérée, non traçable. Ce chapitre pose les conditions de la recevabilité et
les principes qui la garantissent.

\medskip
\textbf{Plan du chapitre.}
\begin{enumerate}
    \item Définir la recevabilité (section~\ref{sec:c06-definir}).
    \item Les conditions (section~\ref{sec:c06-conditions}).
    \item Les principes de l'ACPO (section~\ref{sec:c06-acpo}).
    \item Ce qui rend irrecevable (section~\ref{sec:c06-irrecevable}).
\end{enumerate}

\separateur

% ============================================================================
\section{Définir la Recevabilité}
\label{sec:c06-definir}
% ============================================================================

\begin{boxdef}
\termecle{Recevabilité}~: aptitude d'une preuve à être admise et prise en
compte par la justice. Elle tient à la \emph{manière} dont la preuve a été
obtenue, préservée et documentée, non à sa seule exactitude.
\end{boxdef}

% ============================================================================
\section{Les Conditions}
\label{sec:c06-conditions}
% ============================================================================

\begin{tableaucours}{L{3.2cm}L{9.0cm}}{%
    \tableauHeader{Condition} &
    \tableauHeader{Exigence}
}
Licéité & Obtenue avec un titre (réquisition, mandat, autorisation) \\
Intégrité & Non altérée depuis l'acquisition (condensat) \\
Authenticité & Bien ce qu'elle prétend être, de la source déclarée \\
Traçabilité & Chaîne de possession continue, actes documentés \\
\end{tableaucours}
\vspace{-0.8em}
\begin{center}{\small\itshape Tableau~6.1. Conditions de la recevabilité}\end{center}

% ============================================================================
\section{Les Principes de l'ACPO}
\label{sec:c06-acpo}
% ============================================================================

\begin{boxprocedure}
\textbf{Les quatre principes de l'ACPO}
\begin{enumerate}
    \item Ne pas modifier les données originales.
    \item En cas d'accès nécessaire à l'original, être compétent et expliquer
          l'incidence de ses actes.
    \item Tracer tous les actes~: un tiers doit pouvoir les reproduire et
          obtenir le même résultat.
    \item Le responsable de l'enquête garantit le respect du droit et de ces
          principes.
\end{enumerate}
\end{boxprocedure}

% ============================================================================
\section{Ce qui Rend Irrecevable}
\label{sec:c06-irrecevable}
% ============================================================================

\begin{boxCRO}
\textbf{La recevabilité est l'opposabilité.} La dimension O du Trilemme est la
recevabilité~; la F (fiabilité) recouvre intégrité et authenticité~; la C
(confidentialité) impose la proportionnalité. Une preuve bien conduite vise
\cro{0}{+}{+}.
\end{boxCRO}

\begin{boxscenario}{Opération TSANAGA. Deux pièces, deux sorts}
\textbf{Constatation.} L'ordinateur, acquis selon l'ISO/IEC 27037, à la
chaîne continue, est recevable. La clé USB, dont la chaîne présente un trou
(Chapitre~\ref{chap:custody}), est fragilisée.

\medskip
\textbf{Analyse} (confiance élevée). Les faits tirés de l'ordinateur tiendront~;
ceux tirés de la clé seule seront discutés. On construit la démonstration sur
les pièces les plus solides, et l'on recoupe.
\end{boxscenario}

\begin{boxmetier}{Expert judiciaire et avocat}
\textbf{Anticiper la contestation.} L'expert et l'avocat se demandent, pour
chaque pièce~: « comment la défense pourrait-elle la contester~? ». Ils
corrigent en amont ce qui peut l'être, et signalent honnêtement le reste.
\end{boxmetier}

\separateur

\begin{boxresume}
\textbf{Ce qu'il faut retenir de ce chapitre~:}
\begin{itemize}
    \item \textbf{La recevabilité} tient à la manière, pas à la seule
          exactitude.
    \item \textbf{Quatre conditions}~: licéité, intégrité, authenticité,
          traçabilité.
    \item \textbf{Quatre principes de l'ACPO}~: ne pas altérer, être compétent,
          tracer, garantir.
    \item \textbf{L'opposabilité} du CRO \emph{est} la recevabilité.
    \item \textbf{On anticipe} la contestation.
\end{itemize}
\end{boxresume}

\begin{boxinfo}
\textbf{Points de contrôle de l'élève.} Je sais~: définir la recevabilité~;
citer ses conditions~; énoncer les principes de l'ACPO~; relier au CRO~;
reconnaître une preuve irrecevable.
\end{boxinfo}

\bigskip

\begin{boxexo}[title={\ding{45}\quad Exercice~6.1. Définir \niveaubase}]
À quoi tient la recevabilité d'une preuve~?
\end{boxexo}

\begin{boxexo}[title={\ding{45}\quad Exercice~6.2. Conditions \niveaubase}]
Citez les quatre conditions de la recevabilité.
\end{boxexo}

\begin{boxexo}[title={\ding{45}\quad Exercice~6.3. ACPO \niveaubase}]
Énoncez le premier principe de l'ACPO.
\end{boxexo}

\begin{boxexo}[title={\ding{45}\quad Exercice~6.4. Exacte mais écartée \niveaumoyen}]
Comment une preuve exacte peut-elle être écartée~?
\end{boxexo}

\begin{boxexo}[title={\ding{45}\quad Exercice~6.5. Reproductible \niveaumoyen}]
Pourquoi un tiers doit-il pouvoir reproduire l'analyse~?
\end{boxexo}

\begin{boxexo}[title={\ding{45}\quad Exercice~6.6. Deux pièces \niveaumoyen}]
Pourquoi l'ordinateur est-il recevable et la clé USB fragilisée~?
\end{boxexo}

\begin{boxexo}[title={\ding{45}\quad Exercice~6.7. CRO \niveaumoyen}]
Pourquoi dit-on que l'opposabilité est la recevabilité~?
\end{boxexo}

\begin{boxexo}[title={\ding{45}\quad Exercice~6.8. Anticiper \niveaumoyen}]
Comment anticiper la contestation d'une pièce par la défense~?
\end{boxexo}

\begin{boxexo}[title={\ding{45}\quad Exercice~6.9. Liste de contrôle \niveauexpert}]
Rédigez une liste de contrôle de recevabilité pour TSANAGA.
\end{boxexo}

\begin{boxexo}[title={\ding{45}\quad Exercice~6.10. Note de recevabilité \niveauexpert}]
Rédigez la note qui établit la recevabilité de l'ordinateur saisi.
\end{boxexo}

\bigskip

\begin{boxinfo}
\textbf{Pour aller plus loin.}
\begin{itemize}
    \item ISO/IEC 27037~\cite{iso27037}, 27042~\cite{iso27042}~;
          ACPO~\cite{acpo2012}.
\end{itemize}
\end{boxinfo}

\begin{boxinfo}
\textbf{Transition.} La recevabilité se joue dès l'acquisition~: la
Partie~III s'ouvre sur l'identification et la préservation des sources
(Chapitre~\ref{chap:preserver}).
\end{boxinfo}
